\documentclass[11pt]{article}
\usepackage[margin=0.9in]{geometry}
\usepackage{amsmath,amssymb}
\usepackage{booktabs}
\usepackage{array}
\usepackage{enumitem}
\usepackage{hyperref}
\hypersetup{colorlinks=true,linkcolor=blue,urlcolor=blue}
\emergencystretch=4em
\setlist{nosep,leftmargin=1.3em}

\newcommand{\eps}{\varepsilon}
\newcommand{\shape}[1]{\texttt{(#1)}}
\newcommand{\B}[1]{\textbf{#1}}
\newcommand{\LL}{\mathcal{L}}
\newcommand{\sq}[1]{\Big\lVert #1 \Big\rVert^2}
\newcommand{\readit}[1]{\par\smallskip\noindent\textit{Read it.} #1}

\title{How the corrector is trained, and every way we could change it}
\author{poe\_repair\_min}
\date{2026-09-16}

\begin{document}
\maketitle
\vspace{-1.5em}

\noindent\fbox{\parbox{\dimexpr\textwidth-2\fboxsep-2\fboxrule}{%
Section~1 walks one training step with the real tensor shapes. Section~2 is the same step as four
equations. Section~3 says what is frozen. Section~4 lists what we can change. Section~5 writes the
loss out in full for every variation, one per block, with the code change, what to expect and the
tradeoff.}}

\section{One training step, concretely}

\subsection*{What is on disk}
Sampling was run once per (pair, seed) with no adapter, and every denoising step was recorded.
Per step the cache holds five tensors, each \shape{1, 4, 128, 128} in float32 on CPU:

\begin{center}
\small
\begin{tabular}{ll}
\toprule
\texttt{x\_t} & the noisy latent at this step \\
\texttt{eps\_a\_raw} & what the frozen model predicts given \texttt{"a cat"} \\
\texttt{eps\_b\_raw} & \dots given \texttt{"a dog"} \\
\texttt{eps\_uncond} & \dots given the empty prompt \\
\texttt{eps\_j\_raw} & \dots given \texttt{"a cat and a dog"}. This is the target \\
\bottomrule
\end{tabular}
\end{center}

Four channels at $128\times128$, because a $1024^2$ image is a $128^2$ latent. One sample is
$D = 4\times128\times128 = 65{,}536$ numbers. A cell of $50$ steps is about $250$\,MB.

\subsection*{The step}
Take $K$ cached entries. They may come from different steps, seeds and pairs.

\begin{enumerate}
\item \B{Build the input.} Repeat each \texttt{x\_t} three times and stack:
  \shape{3K, 4, 128, 128}. The three copies will be run against three different prompts.
\item \B{Build the conditioning.} Three prompt embeddings, \texttt{"a cat"}, \texttt{"a dog"},
  empty, each \shape{1, 77, 2048}, concatenated to \shape{3, 77, 2048} and tiled to
  \shape{3K, 77, 2048}. Pooled vectors likewise, \shape{3, 1280} to \shape{3K, 1280}. Timesteps
  \shape{3K}, each sample's $t$ repeated three times.
\item \B{One forward pass} through the adapted model. Output \shape{3K, 4, 128, 128}, cast to fp32
  at once because fp16 cancellation destroys the gradient.
\item \B{Split.} Reshape to \shape{K, 3, 4, 128, 128} and take the three slices: the cat branch
  $a$, the dog branch $b$, the empty branch $u$, each \shape{K, 4, 128, 128}.
\item \B{Compose.} Combine the three into one prediction, \shape{K, 4, 128, 128}. This is the only
  place the three branches meet.
\item \B{Compare to the target}, which is built from cache and never from the adapted model.
  Subtract, square, average over the $65{,}536$ numbers to get \shape{K}, then average over the
  batch to get one number.
\end{enumerate}

\section{The same step, as maths}

\begin{align}
  T &= a + b - u
  &&\text{compose, dial off}
  \label{eq:T}\\
  T^{(w)} &= u + w(a-u) + w(b-u) \;=\; w\,T - (w-1)\,u
  &&\text{compose, dial at } w = 7.5
  \label{eq:Tw}\\
  \text{target} &= \bar u + w(\bar\eps_J - \bar u)
  &&\text{the same dial, on cached tensors}
  \label{eq:tgt}\\
  \LL &= \frac{1}{K}\sum_k \frac{1}{D}\,\sq{T^{(w)}_k - \text{target}_k}
  &&\text{the loss}
  \label{eq:L}
\end{align}

\B{The dial is inside the loss, not outside it.} Both sides are dialled. Setting $w=1$ in
\eqref{eq:Tw} gives back \eqref{eq:T}.

\B{The empty branch is read twice.} The right-hand form of \eqref{eq:Tw} shows it: once inside the
composition, once alone at $-(w-1) = -6.5$.

\B{The composition rule is blind to one move, and the dialled loss is not.}
Add the same tensor $v$ to $a$ and to $u$. Then
\begin{equation}
  T \;\mapsto\; T, \qquad\qquad
  T^{(w)} \;\mapsto\; T^{(w)} - (w-1)\,v .
  \label{eq:gauge}
\end{equation}
The undialled composition does not move at all. The dialled one moves by $-(w-1)v = -6.5v$, and
so does the dialled loss, because its target is built from the cached $\bar u$ and does not follow.

That asymmetry is the problem, and it is worse than blindness. \textbf{The dialled loss can lower
itself by moving the empty branch without making the composition any better.} It is a cheap way to
cancel error that buys nothing at $w=1$. Every variation below either removes the move or leaves
it available.

\section{What is frozen}

\begin{center}
\small
\begin{tabular}{ll}
\toprule
frozen & the entire UNet, both text encoders, the VAE \\
trainable & LoRA factors on $210$ cross-attention projections \\
that is & $9{,}912{,}320$ parameters of $2{,}577{,}376{,}004$, or $0.38\%$ \\
\bottomrule
\end{tabular}
\end{center}

The adapter sits on shared weights, so \B{it changes all three branches at once}. There is no
separate cat adapter. ``Freezing a branch'' therefore does not mean freezing weights. It means not
running that branch through the adapted model and reading the cached frozen tensor instead. That
costs nothing: the tensor is already loaded.

\section{What we can change}

\begin{enumerate}
\item \B{The target} in \eqref{eq:tgt}: the model's own joint-prompt prediction, the noise added to
  a real photo, or something else.
\item \B{The dial}: $w = 7.5$ or $w = 1$.
\item \B{Which branches are adapted}: any of $a$, $b$, $u$ can come from cache instead.
\item \B{An extra term} added to \eqref{eq:L}, holding some branch near a reference.
\item \B{The branch prompts}: \texttt{"a cat"} could be \texttt{"a cat and"}.
\end{enumerate}

\section{Every variation, written out}
\label{sec:vars}

\paragraph{The one convention that makes these readable.}
\B{A bar means the tensor came from the cache}, computed once by the frozen model. No bar means it
came out of this step's forward pass, with the adapter attached. So:
\begin{center}
\small
\begin{tabular}{ll@{\qquad}ll}
\toprule
$a$ & adapted cat branch & $\bar a$ & cached cat branch \\
$b$ & adapted dog branch & $\bar b$ & cached dog branch \\
$u$ & adapted empty branch & $\bar u$ & cached empty branch \\
 & & $\bar\eps_J$ & cached joint-prompt prediction, the target \\
\bottomrule
\end{tabular}
\end{center}
Every $\LL$ below is averaged over the batch and over the $65{,}536$ numbers per sample, exactly as
in \eqref{eq:L}. That averaging is dropped from the display to keep the equations short.

\subsection*{V0: what runs today}
\begin{equation}
  \LL_{V0} \;=\; \sq{\;\underbrace{u + w(a-u) + w(b-u)}_{\text{all three adapted}}
  \;-\;\underbrace{\big[\bar u + w(\bar\eps_J - \bar u)\big]}_{\text{all cached}}\;}
\end{equation}
\readit{Three adapted branches on the left, the cached target on the right, the same dial on both.
Note that $\bar u$ appears on the right and $u$ on the left: the target uses the frozen empty
branch, the composition uses the adapted one.}

\noindent\B{Expect} \quad it trains and the loss falls; four arms ran to $30{,}000$ steps.\\
\B{Tradeoff} \quad the move of \eqref{eq:gauge} is available, and the loss has an incentive to take
it: shifting $u$ cancels error at $w = 7.5$ while doing nothing for the composition at $w = 1$.

\subsection*{V0a: fine the empty branch for moving}
\begin{equation}
  \LL_{V0a} \;=\; \LL_{V0} \;+\; \mu\,\sq{\,u - \bar u\,}, \qquad \mu = 10
\end{equation}
\readit{The same loss, plus a penalty on how far the adapted empty branch has moved from the cached
one.}

\noindent\B{Code} \quad already there, \texttt{trainer.py:496--503}. Free, since $\bar u$ is
cached.\\
\B{Expect} \quad much less drift. It ran, and it does.\\
\B{Tradeoff} \quad a fine is not a ban. What is left of the drift is still multiplied by 6.5, and
that residue can be the size of the whole rest of the error.

\subsection*{V1: freeze the empty branch}
\begin{equation}
  \LL_{V1} \;=\; \sq{\;\underbrace{\bar u + w(a-\bar u) + w(b-\bar u)}_{\text{only } a,b \text{ adapted}}
  \;-\;\big[\bar u + w(\bar\eps_J - \bar u)\big]\;}
  \label{eq:V1}
\end{equation}
Expand and the $\bar u$ terms cancel:
\begin{equation}
  \LL_{V1} \;=\; w^2\,\sq{\;\big(a + b - \bar u\big) \;-\; \bar\eps_J\;}
  \label{eq:V1b}
\end{equation}
\readit{Replace $u$ by $\bar u$ everywhere. The dial then factors out completely, leaving $w^2$
times a loss with no dial in it. That $w^2$ is a constant; the learning rate absorbs it. So
\eqref{eq:V1b} is the board's objective, and V1 is the change that makes ours equal to it.}

\noindent\B{Code} \quad one substitution at \texttt{trainer.py:450}, then drop the third row of the
batch: input becomes \shape{2K, 4, 128, 128}. \B{The sampler must also detach the adapter for its
empty-prompt pass}, or training and sampling disagree.\\
\B{Expect} \quad the same composed predictions as V0. Freeing the empty branch buys nothing:
anything it was doing, the two concept branches can do instead. So no capacity is given up and the
move of \eqref{eq:gauge} is removed outright, not merely priced. One adapter then works at every guidance
setting.\\
\B{Tradeoff} \quad none found. Cheaper by a third, exact where V0a is approximate. The one risk is
the sampler mismatch, a mistake this repo has made before.

\subsection*{V2: hold the empty branch at a plurality prompt}
\begin{equation}
  \LL_{V2} \;=\; \LL_{V0} \;+\; \mu\,\sq{\,u \;-\; \bar\eps(\texttt{"two animals"})\,}
\end{equation}
\readit{V0a with one symbol changed on the right of the penalty. Instead of pulling the empty
branch back to ``no prompt at all'', pull it toward ``a picture with two things in it''.}

\noindent\B{Code} \quad cache one extra tensor per step, then repoint \texttt{trainer.py:497}. It
also lifts a guard: the trainer refuses the anchor with rewritten branch prompts today, because the
only cached reference is the empty prompt's (\texttt{train\_pooled.py:315--321}).\\
\B{Expect} \quad the only variation that both pins the empty branch and moves what the
composition divides by, from ``all pictures'' to ``two-object pictures''. In density terms it is
the only one that reaches the clean two-object product.\\
\B{Tradeoff} \quad the frozen prediction for the \emph{string} \texttt{"two animals"} is not the
same as a proper two-object reference, and how close they are is unchecked. The cache also grows
about 20\%.

\subsection*{V3: adapt only the empty branch}
\begin{equation}
  \LL_{V3} \;=\; \sq{\;\underbrace{u + w(\bar a-u) + w(\bar b-u)}_{\text{only } u \text{ adapted}}
  \;-\;\big[\bar u + w(\bar\eps_J - \bar u)\big]\;}
\end{equation}
\readit{The mirror of V1: bars on the concept branches, none on the empty one. The only free thing
in the whole loss is $u$, one tensor per picture, shared by every pair.}

\noindent\B{Code} \quad input drops to \shape{K, 4, 128, 128}. The cheapest arm by far.\\
\B{Expect} \quad mostly failure, for two reasons. One free tensor shared across all pairs can only
remove the part of the correction common to all pairs, and that part is small and early. And
changing only what you divide by amounts to dividing by the probability the picture has two objects
in it, which rewards the fused animal rather than punishing it.\\
\B{Tradeoff} \quad cheap enough to run for the record, not worth expecting anything from.

\subsection*{V4: the board's version, dial off, all three adapted}
\begin{equation}
  \LL_{V4} \;=\; \sq{\;\big(a + b - u\big) \;-\; \bar\eps_J\;}
\end{equation}
\readit{The dial is gone from both sides. Compare with \eqref{eq:V1b}: the only difference is $u$
against $\bar u$.}

\noindent\B{Code} \quad one argument.\\
\B{Expect} \quad with the empty branch also frozen this \emph{is} V1, differing by the constant
$w^2$. With it free, this and V0 are genuinely different objectives and neither implies the other,
because the term they differ by is exactly the drift $u - \bar u$ scaled by $w-1$.\\
\B{Tradeoff} \quad on its own it fixes nothing. V1 is what makes the dial irrelevant, not $w=1$.

\subsection*{V5: the board's version with real photographs}
\begin{equation}
  \LL_{V5} \;=\; \mathbb{E}_{\,x \sim p_{\text{co}},\; z \sim \mathcal{N}(0,I)}\;
  \sq{\;z \;-\; \big(a + b - u\big)\;},
  \qquad x_t = x + \sigma_t z
\end{equation}
\readit{The target is no longer the model. Take a real photograph $x$ holding both animals as
separate objects, add known noise $z$, and ask the composition to predict the noise you added. The
branches are evaluated at that noised photograph, not at a cached latent, so the cache plays no part
at all.}

\noindent\B{Code} \quad needs a dataset loader and a corpus $p_{\text{co}}$. Neither exists here;
the only images in this repository are SDXL output.\\
\B{Expect} \quad blocked, structurally rather than for a missing download. Almost no pair in this
pool has photographs with both animals, and those that do are look-alike pairs. A pair has such
photos because the two things genuinely co-occur, which is also why the model already composes
them. The data is plentiful where it is not needed.\\
\B{Tradeoff} \quad not applicable until a corpus exists.

\subsection*{V6: a better teacher}
\begin{equation}
  \LL_{V6} \;=\; \mathbb{E}_{\,x \sim p_{\text{filtered}},\; z}\;
  \sq{\;z \;-\; \big(a + b - u\big)\;}
\end{equation}
\readit{V5 with one symbol changed underneath the expectation. The corpus is now the model's own
renders that passed a composition filter, instead of photographs.}

\noindent\B{Code} \quad a render-and-filter pass per pair, then the same loss as V5.\\
\B{Expect} \quad the current target is known to be unreliable; one scope already drops training
cells whose joint-prompt target shows the wrong animals. A filtered sample set is right by
construction exactly where the joint prompt is wrong.\\
\B{Tradeoff} \quad the filter is the risk. The validated scorer counts instances and cannot tell a
cat beside a dog from two dogs, so the corpus needs a second head reading which animals are present.
Without it, this trains toward pictures of the same animal twice. It also changes the target, so
comparisons against V0 stop being clean.

\subsection*{V7: protect single-concept behaviour}
\begin{equation}
  \LL_{V7} \;=\; \big(\text{any loss above}\big)
  \;+\; \nu\,\mathbb{E}_{\,x_t \,\sim\, \text{single-prompt runs}}\;
  \sq{\;\eps_\theta(x_t,t,c_1) \;-\; \bar\eps(x_t,t,c_1)\;}
\end{equation}
\readit{An extra penalty holding the adapted cat branch near the frozen cat branch. The states come
from sampling runs with a single prompt, not from pair runs, and that choice is the whole design.}

\noindent\B{Code} \quad needs a small new cache: states from single-prompt runs plus the frozen
branch output there. Then one extra forward per protected concept per step.\\
\B{Expect} \quad matching the target \emph{requires} the concept branches to move, so this is a
price on the correction rather than a free safeguard. At the pair states it would cancel the fit
outright. At single-concept states it should not, but only if the branch moves differently in the
two places.\\
\B{Tradeoff} \quad it may be actively wrong. If the concept branches are moving because they are
learning the two-object conditioning, this pulls them back from the thing we want. Measure first.

\subsection*{On rewriting the branch prompts}
Three arms ran with the concept prompts changed to \texttt{"a cat and"} against a null of
\texttt{"and"}, and to \texttt{"a cat, two animals"} against a null of \texttt{"two animals"}. In
the equations above that changes which string $a$, $b$ and $u$ were conditioned on, and nothing
else. Renaming a prompt renames a free tensor, so all three reach the same set of composed
predictions; they differ in where training starts, not in where it can end. Their fit losses came
out within 20\% of each other, which is what that predicts. Changing the prompts also invalidates
the cached branch tensors, so it is not free.

\section{All of it in one table}

\begin{center}
\small
\begin{tabular}{c l l l c l}
\toprule
\# & composition & target & extra term & batch & status \\
\midrule
V0  & $u + w(a{-}u) + w(b{-}u)$              & $\bar u + w(\bar\eps_J{-}\bar u)$ & --                        & $3K$ & 4 arms run \\
V0a & same                                    & same                              & $\mu\lVert u-\bar u\rVert^2$ & $3K$ & run \\
V1  & $\bar u + w(a{-}\bar u) + w(b{-}\bar u)$& same                              & --                        & $2K$ & \B{not run} \\
V2  & same as V0                              & same                              & $\mu\lVert u-\bar\eps_{\text{pl}}\rVert^2$ & $3K$ & \B{not run} \\
V3  & $u + w(\bar a{-}u) + w(\bar b{-}u)$     & same                              & --                        & $K$  & not run \\
V4  & $a + b - u$                             & $\bar\eps_J$                      & --                        & $3K$ & $=$ V1 if $u$ frozen \\
V5  & $a + b - u$                             & $z$, real photo                   & --                        & $3K$ & blocked \\
V6  & $a + b - u$                             & $z$, filtered render              & --                        & $3K$ & not run \\
V7  & any                                     & any                               & $\nu\lVert \eps_\theta(c_1)-\bar\eps(c_1)\rVert^2$ & $+1$ & not run \\
\bottomrule
\end{tabular}
\end{center}

\section{Four things not to forget}

\begin{enumerate}
\item $a+b-u$ is unchanged by adding the same $v$ to $a$ and to $u$; $T^{(w)}$ changes by
  $-(w-1)v$. So the undialled losses cannot see that move and the dialled ones can \emph{use} it,
  to cancel error without improving the composition.
\item Freeing the empty branch buys no predictions. It only decides how the answer is split up.
\item Whatever is frozen in training must be frozen in sampling.
\item The branch prompts fix the cache. Change them and the cached branch tensors are wrong.
\end{enumerate}

\section{What to do, in order}

\begin{enumerate}
\item Log $\lVert (a+b-u) - \bar\eps_J\rVert^2$ next to the loss. One line, no forward pass. It says
  whether the loss took the cheap move of \eqref{eq:gauge} instead of fixing the composition.
\item \B{V1.} Cheaper, exact, and it removes the dial from the problem.
\item Two reads against adapters already on disk, neither needing training: does what the adapter
  added to the cat branch point the same way as what the word \texttt{"and"} adds, and is that
  addition the same thing for every concept.
\item \B{V2}, compared against V1.
\item \B{V7}, only if read 3 says the concept branches move differently at pair states and at
  single-concept states.
\item \B{V6} last, because it changes the target.
\end{enumerate}

\paragraph{The numbers behind the claims.}
Deliberately not in here. Two files in this folder carry them.
\texttt{guided-and-unguided-objectives.tex} has the measured drift per arm and what it does to the
loss, the held-out fit, the corpus count and the off-policy probe.
\texttt{correction-loss-variants.tex} has the derivations: why freeing the empty branch buys
nothing, and why renaming prompts changes nothing.

\end{document}
